% LTeX: language=de-DE
% ============================================================
%  Mein Protokoll – nur diese Datei bearbeiten!
%  protokoll.cls muss installiert sein (~/texmf/tex/latex/protokoll/)
% ============================================================
\documentclass{protokoll}
\addbibresource{refs.bib}

% ── Metadaten ────────────────────────────────────────────────
\ueNum{1}
\group{12}
\ueTitle{Bestimmung eines elektrischen Widerstandes}
\supervisor{Dr. Schultze-Bernhardt}
\ueDate{09.03.2026}
\semester{SS 2026}

\authorOne{Maximilian}{Lamprecht}{12129644}
\authorTwo{Manuel}{Lazian}{12524036}   % weglassen bei Einzelprotokoll

% ── Dokument ─────────────────────────────────────────────────
\begin{document}
\maketitle

\section{Aufgabenstellung}

In diesem Versuch soll der elektrische Widerstand eines zur Verfügung
gestellten Resistors, mittels eines Multimeters. Anschließend wird das
Messergebnis mit den Angaben des Herstellers verglichen.

\section{Voraussetzungen und Grundlagen}

Der Farbcode des Resistors muss ausgewertet werden, um die Herstellerangaben zu
Widerstand und Unsicherheiten herauszufinden. Dabei wurde der Farbcode des
Resistors wie folgt von rechts nach links abgelesen (siehe Abbildung
\ref{fig:resistor1}):

\begin{itemize}
    \item Braun
    \item Schwarz
    \item Schwarz
    \item Violet
    \item Gelb
\end{itemize}

\begin{figure}[H]
    \includegraphics[width=0.7\linewidth]{figs/resistor}
    \centering
    \caption{Resistor eingespannt in Bananenstecker. (Gelb links)}
    \label{fig:resistor1}
\end{figure}

In diesem Fall hat der Resistor laut Hersteller einen Wiederstand von
\qty{470}{\ohm} mit einer relativen unsiherheit von \qty{+-1}{\percent}
angegeben. Eine Temperaturkoeffizient ist jedoch nicht angegeben. \\

Dieser Wert wurde mithilfe der Farbcodetabelle aus Abbildung
\ref{fig:color_codes} ermittelt.

\begin{figure}[H]
    \includegraphics[width=0.7\linewidth]{figs/color_codes.png}
    \centering
    \caption{Tabelle für Farbcodes von Wiederständen \cite{res_color_codes}}
    \label{fig:color_codes}
\end{figure}

Für die Auswertung am Ende wird folgende Formel zur Berechnung der Typ B
Unsicherheit benötigt:

\begin{equation}
    \Delta I = I\cdot \mathrm{x}\,\unit{\percent} + \mathrm{y}\,\unit{digits}
    \label{eq:res_error}
\end{equation}

Wobei $\Delta I$ den absoluten Messfehler ist und $I$ sich auf den gemessenen
Wert bezieht. $x$ ist der relative Fehler und $y$ ist die anzahl an Stellen
welche, laut Hersteller, zum absoluten Fehler dazugerechnet werden. In diesem
Fall misst das Gerät maximal eine Nachkommastelle.

Somit ergiebt sich für das Messergebnis $\bar{I}$ mit expliziter absoluter
Unsicherheit folgende Formel:

\begin{equation}
    \bar{I} = I \pm \Delta I
    \label{eq:res_ex_abs}
\end{equation}

\section{Beschreibung der Versuchsanordnung}

Die Enden des Resistors werden jeweils in die Schraublöcher der Bananenstecker
eingespannt. Die Bananenstecker werden dann in die Löcher des Multimeters
(Abbildung \ref{fig:multimeter1}), beschriftet mit \textbf{COM} und
\textbf{V\, \textohm\, \textdegree C}, gesteckt.

\section{Geräteliste}

\begin{figure}[H]
    \centering
    \includegraphics[height=300pt]{figs/multimeter1}
    \caption{Mulimeter}
    \label{fig:multimeter1}
\end{figure}

\begin{table}[H]
    \centering
    \begin{tabular}{| l | l | l | l | l | l | l |}
        \hline
        {Gerät} & {Hersteller} & {Modell} & {InventarNr.} & {Resolution} &
        {Unsicherheit} \\
        \hline\hline
        Multimeter & TE.ELECTRONIC & VA18B & 0109708 & \qty{0.1}{\ohm} &
        $\pm\,($\qty{0.5}{\percent} + \qty{1}{digit(s)}$)$ \\
        \hline
    \end{tabular}
    \caption{Im Versuch verwendete Geräte und Utensilien.}
    \label{tab:devices}
\end{table}

\section{Versuchsdurchführung und Messergebnisse}

Zur Durchführung des Versuches wird das Multimeter auf \textbf{\textohm}, zur
Wiederstandsmessung eingestellt. Dann wird der die digitale Anzeige beobachtet,
bis sich der Wert beruight hat, sodass ein Ergebnis abgelesen wird. Da der
Wert trotzdem im \textit{digit} Bereich fluktuiert, wird der Mittelwert von
Minimum und Maximum genommen und die die Hälfte der Differenz als absolute
Unsicherheit. In diesem Fall ist der Wert zufälligerweise gleich dem Modus.

Dieser Prozess wird dann wiederholt, wobei die Einspannung des Resistors
umgekehrt in das Multimeter gesteckt wird, um zu validieren, dass der
Wiederstand in beide Stromflussrichtungen gleich startk ist. Wie in Tabelle
\ref{tab:measurements} zu sehen, ist dies der Fall.

\begin{table}[H]
    \centering
    \begin{tabular}{| l | S | S |}
        \hline
        {Messung} & {Messwert / \unit\ohm} & {Unsicherheit / \unit\ohm} \\
        \hline\hline
        Wiederstand (Braun\, \textrightarrow\, COM) & 468.5 & +-0.1 \\
        \hline
        Wiederstand (Gelb\, \textrightarrow\, COM) & 468.5 & +-0.1 \\
        \hline
    \end{tabular}
    \caption{Messergebnisse}
    \label{tab:measurements}
\end{table}

\section{Auswertung}

Da das Multimeter bereits die gewünschte Größe, in diesem Fall den Wiederstand,
misst, muss nur noch die Typ B unsicherheit ausgerechnet werden. Dazu wird der
Messwert $I$ erst in Formel \ref{eq:res_error} eingesetzt und $\mathrm{x}$ und
$\mathrm{y}$ werden den Angaben des Herstellers des Multimeters nach angepasst.
In diesem Fall sind, bis \qty{60}{\Mohm}, $\mathrm{x}=\,$\qty{0.5}{\percent}
und $\mathrm{y}=\,$\qty{1}{\unit{digit}}. Daraus ergibt sich mit einer Auflösung
von \qty{0.1}{\ohm} und dem Messwert $I$ durch einsetzen der in Formel
\ref{eq:res_error} für den absoluten Fehler:

\begin{equation*}
    \Delta I = (\num{2.2325} + \num{0.1})\,\unit\ohm \approx \qty{2.3}{\ohm}
\end{equation*}

Da die Unsicherheit durch Fluktuationen viel geringer ist als die Typ B
Unsicherheit, wird diese vernachlässigt. Des weiteren, da beide Messwerte
inklusive deren Fluktuationsunsicherheiten identisch sind, wird der finale Wert
$\bar{I}$ nur einmal ausgerechnet.

Somit lautet das finale Ergebnis der Messung:

\begin{equation*}
    \uuline{\bar{I} = \qty{468.5(2.3)}{\ohm}}
\end{equation*}

\section{Diskussion}

\begin{table}[H]
    \centering
    \begin{tabular}{| l | S | S |}
        \hline
        Größe & {Herstellerwert / \unit\ohm} & {Messergebnis / \unit\ohm} \\
        \hline\hline
        $I$ & 470(5) & 468.5(2.3) \\
        \hline
    \end{tabular}
    \caption{Vergleich: Herstellerangaben vs Messergebnis}
    \label{tab:comparison}
\end{table}

Die Angaben des Herstellers des Resistors besagen, dass dieser eine maximale
Unsicherheit von \qty{+-1}{\percent} aufweist. Bei \qty{470}{\ohm} beträgt diese
\qty{+-4.7}{\ohm}. Damit liegt das Messergebnis mit \qty{468.5(2.4)}{\ohm}
innerhalb des vom Hersteller angegebenen Konfidenzintervalls.

\section{Zusammenfassung}

In diesem Experiment wurde der elektrische Wiederstand eines Resistors gemessen.
Der Resistor is mit \qty{470}{\ohm} \qty{+-1}{\percent} bzw. \qty{470(5)}{\ohm}.
Das Messergebnis liegt somit mit \qty{468.5(2.3)}{\ohm} innerhalb der Angaben
des Herstellers.

\pagebreak
\section{Literatur}

\nocite{*}
\printbibliography

\end{document}
